% ============================================================
%  RIJEKA — Quantitative Methodology & System Documentation
%  v1.1 | Sprint 1 Complete
%  LaTeX source — compile with: pdflatex (twice for TOC/refs)
% ============================================================
\documentclass[11pt, a4paper]{article}

% ── Packages ────────────────────────────────────────────────
\usepackage[T1]{fontenc}
\usepackage[utf8]{inputenc}
\usepackage[margin=2.5cm, top=3cm, bottom=3cm, headheight=15pt]{geometry}
\usepackage{amsmath, amssymb, mathtools, bm}
\usepackage{xcolor}
\usepackage{booktabs, array, longtable, multirow}
\usepackage{listings}
\usepackage{fancyhdr}
\usepackage{titlesec}
\usepackage{tcolorbox}
\usepackage{mdframed}
\usepackage{enumitem}
\usepackage{hyperref}
\usepackage{microtype}
\usepackage{parskip}
\usepackage{graphicx}
\usepackage{float}

\tcbuselibrary{skins, breakable}

% ── Colour palette ──────────────────────────────────────────
\definecolor{navy}{HTML}{0D2137}
\definecolor{teal}{HTML}{0A7A5E}
\definecolor{lteal}{HTML}{E6F7F3}
\definecolor{amber}{HTML}{C07800}
\definecolor{ltamber}{HTML}{FFF8E6}
\definecolor{rblue}{HTML}{1E5A9E}
\definecolor{ltblue}{HTML}{EEF4FF}
\definecolor{dgrey}{HTML}{5A6A7A}
\definecolor{lgrey}{HTML}{F5F7FA}
\definecolor{mgrey}{HTML}{D8DFE8}
\definecolor{codebg}{HTML}{F0F5F0}
\definecolor{codegreen}{HTML}{1A5C28}
\definecolor{red}{HTML}{B03020}

% ── Hyperref ────────────────────────────────────────────────
\hypersetup{
  colorlinks=true,
  linkcolor=navy,
  urlcolor=teal,
  citecolor=teal,
  pdftitle={Rijeka Quantitative Methodology v1.1},
  pdfauthor={Rijeka Core Team},
  bookmarksnumbered=true
}

% ── Fonts ───────────────────────────────────────────────────
\usepackage{charter}

% ── Title formatting ────────────────────────────────────────
\titleformat{\section}
  {\color{navy}\large\bfseries}
  {\thesection}{1em}{}
  [\color{teal}\titlerule]

\titleformat{\subsection}
  {\color{navy}\normalsize\bfseries}
  {\thesubsection}{1em}{}

\titleformat{\subsubsection}
  {\color{teal}\normalsize\bfseries}
  {\thesubsubsection}{1em}{}

\titlespacing*{\section}{0pt}{18pt}{6pt}
\titlespacing*{\subsection}{0pt}{12pt}{4pt}
\titlespacing*{\subsubsection}{0pt}{10pt}{3pt}

% ── Header / Footer ─────────────────────────────────────────
\pagestyle{fancy}
\fancyhf{}
\fancyhead[L]{\small\color{dgrey}\textbf{RIJEKA} \; | \; Quantitative Methodology \& System Documentation}
\fancyhead[R]{\small\color{dgrey}v1.1 \; | \; Sprint 1}
\fancyfoot[L]{\small\color{dgrey}OPEN SOURCE \; | \; rijeka.io}
\fancyfoot[R]{\small\color{dgrey}\thepage}
\renewcommand{\headrulewidth}{0.4pt}
\renewcommand{\headrule}{\hbox to\headwidth{\color{teal}\leaders\hrule height \headrulewidth\hfill}}
\renewcommand{\footrulewidth}{0pt}

% ── Code listings ───────────────────────────────────────────
\lstset{
  backgroundcolor=\color{codebg},
  basicstyle=\footnotesize\ttfamily\color{codegreen},
  breaklines=true,
  frame=single,
  framerule=0.4pt,
  rulecolor=\color{mgrey},
  xleftmargin=6pt,
  xrightmargin=6pt,
  aboveskip=6pt,
  belowskip=6pt,
  showstringspaces=false,
  keywordstyle=\color{rblue}\bfseries,
  commentstyle=\color{dgrey}\itshape,
  stringstyle=\color{teal},
  language=Python
}

% ── Custom boxes ────────────────────────────────────────────
\tcbset{
  definition/.style={
    colback=lteal, colframe=teal, coltext=navy,
    fonttitle=\bfseries\small, title={Definition},
    arc=2pt, boxrule=0.6pt, left=6pt, right=6pt,
    top=4pt, bottom=4pt, breakable
  },
  important/.style={
    colback=ltblue, colframe=rblue, coltext=navy,
    fonttitle=\bfseries\small,
    arc=2pt, boxrule=0.6pt, left=6pt, right=6pt,
    top=4pt, bottom=4pt, breakable
  },
  warning/.style={
    colback=ltamber, colframe=amber, coltext=navy,
    fonttitle=\bfseries\small, title={Note},
    arc=2pt, boxrule=0.6pt, left=6pt, right=6pt,
    top=4pt, bottom=4pt, breakable
  },
  formula/.style={
    colback=lgrey, colframe=mgrey, coltext=navy,
    arc=2pt, boxrule=0.6pt, left=10pt, right=10pt,
    top=6pt, bottom=6pt, breakable
  },
  sprintbox/.style={
    colback=lgrey, colframe=mgrey,
    arc=2pt, boxrule=0.6pt, left=6pt, right=6pt,
    top=3pt, bottom=3pt
  }
}

% ── Math operators ───────────────────────────────────────────
% Math sets
\newcommand{\Prob}{\mathbb{P}}
\newcommand{\dd}{\,\mathrm{d}}
\newcommand{\T}{\mathrm{T}}

% ── Table column types ───────────────────────────────────────
\newcolumntype{L}[1]{>{\raggedright\arraybackslash}p{#1}}
\newcolumntype{C}[1]{>{\centering\arraybackslash}p{#1}}
\newcolumntype{R}[1]{>{\raggedleft\arraybackslash}p{#1}}

% ── Sprint status command ────────────────────────────────────
\newcommand{\sprintstatus}[3]{%
  \begin{tcolorbox}[sprintbox]
  \small
  \textbf{\color{navy}Sprint #1} \quad
  \textbf{\color{#2}#3}
  \end{tcolorbox}\vspace{4pt}
}

% ============================================================
\begin{document}
% ============================================================

% ── Cover page ───────────────────────────────────────────────
\thispagestyle{empty}
\vspace*{3cm}
\begin{center}
  {\Huge\bfseries\color{navy} RIJEKA}\\[0.4cm]
  {\large\color{teal} Open-Source Full Revaluation Risk System}\\[0.2cm]
  \textcolor{mgrey}{\rule{12cm}{0.8pt}}\\[1cm]
  {\LARGE\bfseries\color{navy} Quantitative Methodology \&}\\[0.3cm]
  {\LARGE\bfseries\color{navy} System Documentation}\\[2cm]
  \begin{center}
  \begin{tabular}{ll}
    \toprule
    \textbf{Document Version} & v1.1 \\
    \textbf{Status}           & Active — updated each sprint \\
    \textbf{Sprint Coverage}  & Sprint 1 complete; Sprints 2–20 framework defined \\
    \textbf{Reference Date}   & 2025-01-15 \\
    \textbf{ISDA SIMM}        & v2.6 \\
    \textbf{QuantLib}         & 1.34 \\
    \textbf{Classification}   & Open Source — Public \\
    \textbf{Primary Audience} & Model validators, quants, risk managers \\
    \bottomrule
  \end{tabular}
  \end{center}
  \vspace{1.5cm}
  \textit{\color{dgrey}This document is the single source of truth for all quantitative\\
  methodology, model validation support, and data lineage in Rijeka.\\[0.2cm]
  \textbf{\color{teal}A feature is not production-ready until its methodology section is written.}}
\end{center}

\newpage
\tableofcontents
\newpage

% ============================================================
\section{Purpose and Scope}
% ============================================================

\subsection{Document Purpose}

This document simultaneously serves five audiences:

\begin{center}
\begin{tabular}{L{4.5cm} L{4.5cm} L{4cm}}
\toprule
\textbf{Purpose} & \textbf{Audience} & \textbf{Standard} \\
\midrule
Quantitative methodology & Quants, model validators & Industry best practice \\
Model validation support & MRM teams, counterparties & SR 11-7, SS1/23 \\
System configuration ref & Risk managers, technology & Internal governance \\
Data lineage documentation & Auditors, regulators & BCBS 239 \\
Regulatory alignment & Fed, ECB, FCA, CFTC & BCBS/IOSCO, EMIR \\
\bottomrule
\end{tabular}
\end{center}

\subsection{ISDA SIMM as Universal Vocabulary}

Rijeka uses ISDA SIMM v2.6 naming conventions throughout the entire system.
Every risk class, sensitivity type, tenor, and sub-curve follows SIMM definitions,
ensuring a model validator familiar with ISDA SIMM can immediately orient themselves
without learning a parallel vocabulary.

\begin{tcolorbox}[important, title={SIMM Sensitivity Notation}]
\begin{align}
s_{ik}     &\quad \text{Raw sensitivity of instrument } i \text{ to risk factor } k \\
WS_{ik}    &= RW_k \cdot s_{ik} \quad \text{(weighted sensitivity)} \\
K_b        &= \sqrt{\sum_i WS_{ib}^2 + \sum_{i \neq j} \rho_{ij}\, WS_{ib}\, WS_{jb}}
             \quad \text{(bucket margin)} \\
IM_{rc}    &= \sqrt{\sum_b K_b^2 + \sum_{b \neq c} \gamma_{bc}\, S_b\, S_c}
             \quad \text{(risk class margin)} \\
IM_{total} &= \sum_{rc} IM_{rc} \quad \text{(total SIMM IM)}
\end{align}
where $\rho_{ij}$ is the within-bucket correlation and $\gamma_{bc}$ is the
cross-bucket correlation from the SIMM v2.6 parameter set.
\end{tcolorbox}

% ============================================================
\section{System Architecture Overview}
% ============================================================

\subsection{Dual-Mode Computation Engine}

Every analytics module supports two computation modes, selected at runtime.

\subsubsection{Mode 1: Full Revaluation}

For each trade $i$ and scenario $\omega$:
\begin{equation}
  \Delta P_i(\omega) = \mathrm{NPV}_i\!\left(\mathcal{M}(\omega)\right)
                     - \mathrm{NPV}_i\!\left(\mathcal{M}_0\right)
\end{equation}
where $\mathcal{M}(\omega)$ is the full market data state under scenario $\omega$
and $\mathcal{M}_0$ is the base market data. Exact — captures all nonlinearity.
Used for EOD, regulatory submissions, LCH IM, and backtesting.

\subsubsection{Mode 2: Sensitivities (Taylor Approximation)}

\begin{equation}
  \Delta P \approx
    \underbrace{\sum_i s_i \,\Delta x_i}_{\text{delta}} +
    \underbrace{\frac{1}{2}\sum_{i,j} \Gamma_{ij}\,\Delta x_i\,\Delta x_j}_{\text{gamma}} +
    \underbrace{\mathcal{V}\,\Delta\sigma}_{\text{vega}} +
    \underbrace{\Theta\,\Delta t}_{\text{theta}}
\end{equation}

Scenarios $\Delta \mathbf{x}$ are drawn from the SIMM covariance matrix:
\begin{equation}
  \Sigma_{ij} = \sigma_i\,\sigma_j\,\rho_{ij}
\end{equation}
where $\sigma_i = RW_i$ (SIMM risk weight, acting as volatility proxy)
and $\rho_{ij}$ is the SIMM within-bucket correlation. Used for intraday
monitoring, what-if analysis, and limit checking.

\subsection{Multi-Curve Framework}

Post-2008, the discount curve and projection curve are strictly separate objects.

\begin{tcolorbox}[definition]
\textbf{Discount curve:} the OIS curve for the CSA collateral currency.
\begin{equation}
  D(t) = \exp\!\left(-r_{\mathrm{OIS}}(t)\cdot t\right)
\end{equation}
\textbf{Projection curve:} the index-specific forward curve.
\begin{equation}
  F(t_1, t_2) = \frac{1}{\tau(t_1,t_2)}
                \left(\frac{D(t_1)}{D(t_2)} - 1\right)
\end{equation}
where $\tau(t_1,t_2)$ is the day count fraction under the relevant convention.
\end{tcolorbox}

The NPV of a floating leg in currency $c$:
\begin{equation}
  \mathrm{NPV}_{\mathrm{float}} =
    N \sum_{i=1}^{n} \tau_i\, F(t_{i-1},\, t_i)\, D(t_i)
\end{equation}
where $N$ is notional, $D(t_i)$ comes from the OIS curve (CSA currency),
and $F(t_{i-1}, t_i)$ comes from the index projection curve
(e.g.\ USD\_TSOFR\_3M for a Term SOFR 3M IRS).

% ============================================================
\section{Market Data Architecture}
% ============================================================

\subsection{Three-Store Design}

\begin{center}
\begin{tabular}{L{2.4cm} L{3.5cm} L{2.2cm} L{2.2cm} L{2.5cm}}
\toprule
\textbf{Store} & \textbf{Purpose} & \textbf{Horizon} & \textbf{Mutability} & \textbf{Table} \\
\midrule
History    & VaR, stress, LCH FHS, backtesting
           & Up to 10y (2520 bd)
           & \textbf{Immutable} — DB trigger blocks UPDATE/DELETE
           & \texttt{market\_history} \\
\addlinespace
Production & Committed EOD/intraday snapshots for all analytics
           & Rolling 90-day hot cache
           & Immutable once committed; amendment = new version
           & \texttt{production\_snapshots} \\
\addlinespace
Working    & Live quotes, staging, what-if
           & Session only
           & Freely mutable; truncated each morning
           & \texttt{working\_quotes} \\
\bottomrule
\end{tabular}
\end{center}

EOD commit is \emph{atomic}: writes Production and History simultaneously.
Both succeed or both roll back. The stores can never be out of sync.

\subsection{Data Quality Tier System}

Every data point carries an immutable quality tier stamped at ingestion.
Risk engines enforce minimum tier requirements before any run.

\begin{center}
\begin{tabular}{C{0.5cm} L{3.5cm} L{5.5cm} L{2.8cm}}
\toprule
\textbf{T} & \textbf{Name} & \textbf{Sources} & \textbf{Regulatory VaR} \\
\midrule
1 & LIVE\_VENDOR & Bloomberg BDP/BDH, Refinitiv & Permitted \\
2 & OFFICIAL\_CB & FRED, ECB, BOE, RBA, SNB, MAS & Permitted \\
3 & VENDOR\_IMPORT & Bloomberg CSV, Refinitiv CSV & Permitted \\
4 & INTERNAL\_IMPORT & Internal systems, broker axes & Permitted (with approval) \\
5 & SYNTHETIC\_SEED & Rijeka-generated (testing only) & \textbf{Blocked} \\
6 & GAP\_FILLED & Carry-forward / interpolation & \textbf{Blocked} \\
\bottomrule
\end{tabular}
\end{center}

\subsection{VaR Readiness Check}

Before any historical simulation, the system evaluates:
\begin{equation}
  \text{coverage} = \frac{|\mathcal{D}_{\text{available}} \cap \mathcal{D}_{\text{required}}|}{|\mathcal{D}_{\text{required}}|}
\end{equation}
where $\mathcal{D}_{\text{required}}$ is the set of required business dates
(typically 1,250 days for VaR or 2,520 for LCH FHS) and
$\mathcal{D}_{\text{available}}$ is the set of dates with at least Tier 1–4
data for all required risk factors. Regulatory runs are blocked if
$\text{coverage} < 1.00$; indicative runs warn if $\text{coverage} < 0.99$.

% ============================================================
\section{Sprint 1: Market Data Workspace}
% ============================================================

\sprintstatus{1}{teal!70!black}{COMPLETE — OIS curves, surfaces, fixings, snap, snapshots}

\subsection{Rates Curve Universe}

\begin{center}
\begin{tabular}{L{2.8cm} R{1.2cm} L{8cm}}
\toprule
\textbf{Class} & \textbf{Count} & \textbf{Description} \\
\midrule
OIS (INDEPENDENT) & 13 & DM overnight indexed swap mother curves \\
BASIS              & 22 & Single-currency IBOR/term RFR vs OIS \\
XCCY\_BASIS        & 12 & Cross-currency basis — all DM pairs vs USD \\
FUNDING            &  7 & Firm cost-of-funds: OIS + credit spread \\
\midrule
\textbf{Total}     & \textbf{54} & \\
\bottomrule
\end{tabular}
\end{center}

Bootstrap execution order follows topological sort of the dependency graph:
\begin{equation}
  \underbrace{\text{Pass 1: OIS}}_{\text{independent}}
  \;\to\;
  \underbrace{\text{Pass 2: Basis}}_{\text{needs OIS}}
  \;\to\;
  \underbrace{\text{Pass 3: XCCY}}_{\text{needs OIS + FX}}
  \;\to\;
  \underbrace{\text{Pass 4: Funding}}_{\text{needs OIS, no solver}}
\end{equation}

% ----------------------------------------------------------
\subsection{OIS Curve Framework}
% ----------------------------------------------------------

\subsubsection{Discount Factor and Zero Rate}

\begin{tcolorbox}[definition, title={Fundamental Curve Relations}]
Let $D(t)$ denote the discount factor for maturity $t$ (in years).
The continuously compounded zero rate is:
\begin{equation}
  r(t) = -\frac{\ln D(t)}{t}
  \qquad\Longleftrightarrow\qquad
  D(t) = e^{-r(t)\,t}
\end{equation}
The instantaneous forward rate:
\begin{equation}
  f(t) = -\frac{\partial}{\partial t}\ln D(t)
       = r(t) + t\,\frac{\partial r}{\partial t}
\end{equation}
The simply compounded period forward rate for $[t_1, t_2]$:
\begin{equation}
  F(t_1, t_2) = \frac{1}{\tau(t_1,t_2)}
                \left(\frac{D(t_1)}{D(t_2)} - 1\right)
\end{equation}
The par OIS swap rate for maturity $T$ with fixed payment times $\{t_i\}$:
\begin{equation}
  S(T) = \frac{D(0) - D(T)}{\sum_{i=1}^{n} \tau_i\, D(t_i)}
       = \frac{1 - D(T)}{A(T)}
\end{equation}
where $A(T) = \sum_{i=1}^{n}\tau_i D(t_i)$ is the annuity (DV01 numeraire).
\end{tcolorbox}

\subsubsection{Bootstrap as Root-Finding Problem}

Given market quotes $\{(T_j,\, Q_j)\}_{j=1}^{n}$ — where $Q_j$ is the observed
par OIS swap rate for maturity $T_j$ — the bootstrap solves for discount factors
$\{D(T_j)\}$ such that each instrument reprices to par:

\begin{equation}
  \mathcal{F}_j\!\left(D(T_1),\ldots,D(T_j)\right)
  = Q_j^{\text{model}} - Q_j^{\text{market}} = 0
  \qquad j = 1,\ldots,n
\end{equation}

For an OIS swap this gives an \emph{exact} closed-form solution for $D(T_j)$
given $D(T_1),\ldots,D(T_{j-1})$ (sequential bootstrap, no iteration needed
for piecewise-constant forwards). QuantLib uses Brent's method for general
interpolation methods where no closed form exists:

\begin{equation}
  D(T_j) = \arg\min_{D} \left|S^{\text{model}}(T_j;\, D) - Q_j\right|
\end{equation}

with convergence tolerance $\varepsilon = 10^{-12}$ on the rate residual.

\subsubsection{OIS Curve Conventions}

\begin{center}
\small
\begin{tabular}{L{1.1cm} L{1.3cm} L{3.2cm} L{2.0cm} L{2.8cm} C{0.8cm} C{0.8cm} C{1.0cm}}
\toprule
\textbf{Ccy} & \textbf{Index} & \textbf{QL Class} & \textbf{Day Count} & \textbf{Calendar}
& \textbf{T+} & \textbf{Lag} & \textbf{Tele.} \\
\midrule
USD & SOFR   & \texttt{ql.Sofr()}  & Act/360   & US(FedRes)     & 1 & 2 & Yes \\
EUR & €STR   & \texttt{ql.Estr()}  & Act/360   & TARGET()       & 2 & 2 & Yes \\
GBP & SONIA  & \texttt{ql.Sonia()} & Act/365F  & UK(Exchange)   & 0 & 0 & No  \\
JPY & TONAR  & \texttt{ql.Tonar()} & Act/365F  & Japan()        & 2 & 2 & Yes \\
CHF & SARON  & \texttt{ql.Saron()} & Act/360   & Switzerland()  & 2 & 2 & Yes \\
AUD & AONIA  & \texttt{ql.Aonia()} & Act/365F  & Australia()    & 0 & 0 & Yes \\
CAD & CORRA  & \texttt{ql.Corra()} & Act/365F  & Canada()       & 1 & 2 & Yes \\
SEK & SWESTR & \texttt{OvernightIndex} & Act/360 & Sweden()     & 1 & 1 & Yes \\
NOK & NOWA   & \texttt{OvernightIndex} & Act/360 & Norway()     & 1 & 1 & Yes \\
DKK & DESTR  & \texttt{OvernightIndex} & Act/360 & Denmark()    & 1 & 1 & Yes \\
NZD & NZIONA & \texttt{OvernightIndex} & Act/365F & NZ()        & 0 & 0 & Yes \\
HKD & HONIA  & \texttt{OvernightIndex} & Act/365F & HongKong()  & 0 & 0 & Yes \\
SGD & SORA   & \texttt{OvernightIndex} & Act/365F & Singapore() & 0 & 0 & Yes \\
\bottomrule
\end{tabular}
\end{center}

\begin{tcolorbox}[warning, title={GBP SONIA Convention Note}]
GBP SONIA is unique among the DM OIS curves: settlement is T+0 and payment lag is 0 days.
This reflects Bank of England convention. In QuantLib, set
\texttt{telescopicValueDates=False} for SONIA — all other DM OIS curves use \texttt{True}.
\end{tcolorbox}

% ----------------------------------------------------------
\subsection{Interpolation Methods — Mathematical Specification}
% ----------------------------------------------------------

Eight interpolation methods are supported. The choice determines the functional
form of $D(t)$ (or $r(t)$ or $f(t)$) between pillar dates.
Let $\{t_0 < t_1 < \cdots < t_n\}$ be the bootstrap pillars.

\subsubsection{1. Log-Linear Discount (Default)}

Interpolates $\ln D(t)$ linearly between pillars:
\begin{equation}
  \ln D(t) = \ln D(t_i) + \frac{t - t_i}{t_{i+1} - t_i}
             \left[\ln D(t_{i+1}) - \ln D(t_i)\right],
  \qquad t \in [t_i, t_{i+1}]
\end{equation}
\begin{equation}
  \Longrightarrow \quad f(t) = \frac{\ln D(t_i) - \ln D(t_{i+1})}{t_{i+1} - t_i}
  = \text{const on } [t_i, t_{i+1}]
\end{equation}
Instantaneous forward rates are \emph{piecewise constant} and \emph{positive} by construction.
Industry standard for OIS bootstrapping. QuantLib class: \texttt{PiecewiseLogLinearDiscount}.

\subsubsection{2. Linear Zero}

Interpolates the zero rate $r(t)$ linearly:
\begin{equation}
  r(t) = r(t_i) + \frac{t-t_i}{t_{i+1}-t_i}\left[r(t_{i+1})-r(t_i)\right]
\end{equation}
Simple and fast. Instantaneous forwards are piecewise constant (same structure as
log-linear discount but on rates, not log-DFs). QuantLib: \texttt{PiecewiseLinearZero}.

\subsubsection{3. Cubic Spline Zero}

Interpolates $r(t)$ with a natural cubic spline. The spline satisfies:
\begin{equation}
  r(t) = a_i + b_i(t-t_i) + c_i(t-t_i)^2 + d_i(t-t_i)^3,
  \qquad t \in [t_i, t_{i+1}]
\end{equation}
with $C^2$ continuity: $r$, $r'$, and $r''$ all continuous at each pillar.
Natural boundary conditions: $r''(t_0) = r''(t_n) = 0$.
Produces smooth, continuous forward rates. Risk: Runge phenomenon at endpoints.
QuantLib: \texttt{PiecewiseCubicZero}.

The instantaneous forward rate is:
\begin{equation}
  f(t) = r(t) + t\,r'(t)
       = a_i + (2b_i + a_i/t)(t-t_i) + \cdots
\end{equation}
where $r'(t)$ is the spline derivative — smooth everywhere.

\subsubsection{4. Log-Cubic Discount}

Interpolates $\ln D(t)$ with a cubic spline (same $C^2$ conditions as above
but applied to the log-discount):
\begin{equation}
  \ln D(t) = a_i + b_i(t-t_i) + c_i(t-t_i)^2 + d_i(t-t_i)^3
\end{equation}
\begin{equation}
  f(t) = -\frac{\partial}{\partial t}\ln D(t) = -(b_i + 2c_i(t-t_i) + 3d_i(t-t_i)^2)
\end{equation}
Forwards are smooth but not guaranteed positive. Good balance of smoothness and
stability. QuantLib: \texttt{PiecewiseLogCubicDiscount}.

\subsubsection{5. Monotone Cubic Zero (Hyman Filter)}

Cubic spline on $r(t)$ with Hyman (1983) monotonicity filter applied to spline
coefficients. For each interval, the filter constrains $c_i$ and $d_i$ such that
the interpolant preserves local monotonicity:
\begin{equation}
  \text{if } r(t_i) \leq r(t_{i+1}) \text{ then }
  r(t) \text{ is non-decreasing on } [t_i, t_{i+1}]
\end{equation}
The filter replaces $c_i$ with:
\begin{equation}
  \tilde{c}_i = \min\!\left(3\,\frac{r(t_{i+1})-r(t_i)}{h_i},\; c_i\right)
\end{equation}
Eliminates oscillations without global dampening. QuantLib: \texttt{PiecewiseMonotoneCubicZero}.

\subsubsection{6. Flat Forward}

Interpolates the instantaneous forward rate $f(t)$ as piecewise constant, solving
directly for the forward level that reprices each instrument:
\begin{equation}
  f(t) = f_i = \text{const}, \quad t \in (t_{i-1}, t_i]
\end{equation}
\begin{equation}
  D(t) = D(t_{i-1})\,\exp\!\left(-f_i\,(t - t_{i-1})\right)
\end{equation}
Maximum locality — changing one pillar affects only the adjacent segment.
QuantLib: \texttt{PiecewiseFlatForward}.

\subsubsection{7. Convex Monotone}

Andersen \& Piterbarg (2010) convex monotone interpolation on $f(t)$.
Parameterises the forward curve as:
\begin{equation}
  f(t) = \phi_i + \psi_i\,g\!\left(\frac{t-t_{i-1}}{h_i}\right),
  \quad g(x) = \frac{x^2 - x}{2},
  \quad h_i = t_i - t_{i-1}
\end{equation}
where $\phi_i$, $\psi_i$ are chosen to ensure $f(t) > 0$ everywhere and
convexity is preserved between pillars. QuantLib: \texttt{PiecewiseConvexMonotone}.

\subsubsection{8. Kruger Cubic Zero}

Kruger (2002) locally-monotone cubic interpolation on $r(t)$. Unlike the global
Hyman filter, Kruger's algorithm is \emph{local}: it computes each spline slope
from adjacent data only, avoiding the global dampening effect:
\begin{equation}
  r'(t_i) = \begin{cases}
    \frac{h_{i-1}\,\delta_i + h_i\,\delta_{i-1}}{h_{i-1}+h_i}
    & \text{if } \delta_{i-1}\,\delta_i > 0 \\
    0 & \text{otherwise}
  \end{cases}
\end{equation}
where $\delta_i = (r(t_{i+1})-r(t_i))/h_i$ is the divided difference.
QuantLib: \texttt{PiecewiseKrugerZero}.

\subsubsection{Interpolation Method Comparison}

\begin{center}
\small
\begin{tabular}{L{3.2cm} L{2cm} L{2cm} L{2.2cm} L{3cm}}
\toprule
\textbf{Method} & \textbf{Interpolates} & \textbf{Forward} & \textbf{Positive $f$?} & \textbf{Best For} \\
\midrule
Log-Linear Discount & $\ln D(t)$ & Piecewise const & \textbf{Guaranteed} & OIS — \textbf{Default} \\
Linear Zero & $r(t)$ & Piecewise const & Not guaranteed & Simple structures \\
Cubic Spline Zero & $r(t)$ & Smooth ($C^1$) & Not guaranteed & Smooth fwd profile \\
Log-Cubic Discount & $\ln D(t)$ & Smooth ($C^1$) & Not guaranteed & Gov't curves \\
Monotone Cubic & $r(t)$ & Smooth, monotone & Conditionally & Curves without oscillation \\
Flat Forward & $f(t)$ & Step function & \textbf{Guaranteed} & Max locality \\
Convex Monotone & $f(t)$ & Smooth, convex & \textbf{Guaranteed} & Central bank curves \\
Kruger Cubic & $r(t)$ & Smooth, local & Conditionally & Alt. to Hyman filter \\
\bottomrule
\end{tabular}
\end{center}

% ----------------------------------------------------------
\subsection{Single-Currency Basis Curves}
% ----------------------------------------------------------

\sprintstatus{3}{amber!80!black}{PLANNED — Bootstrap wired Sprint 3}

\subsubsection{Simultaneous Dual-Curve Bootstrap}

The basis curve is not bootstrapped independently. It is solved simultaneously
with the OIS discount curve via the following system. Let $D^{\mathrm{OIS}}(t)$
denote the already-bootstrapped OIS discount factors. The basis projection
curve $D^{\mathrm{basis}}(t)$ is found by solving:

\begin{equation}
  S^{\mathrm{IRS}}(T) =
    \frac{\sum_{i=1}^{n^f}\tau_i^f\, F^{\mathrm{basis}}(t_{i-1}^f, t_i^f)\,
          D^{\mathrm{OIS}}(t_i^f)}
         {\sum_{j=1}^{n^f}\tau_j^f\, D^{\mathrm{OIS}}(t_j^f)}
    = Q_T
\end{equation}

where $F^{\mathrm{basis}}$ is the forward rate implied by $D^{\mathrm{basis}}$:
\begin{equation}
  F^{\mathrm{basis}}(t_1,t_2)
  = \frac{1}{\tau(t_1,t_2)}
    \left(\frac{D^{\mathrm{basis}}(t_1)}{D^{\mathrm{basis}}(t_2)} - 1\right)
\end{equation}

The OIS curve provides discounting; the basis curve provides projection.
They are decoupled — $D^{\mathrm{OIS}}$ enters as a fixed input to the
basis bootstrap.

\subsubsection{Tenor Basis Spreads}

For a tenor basis swap (e.g.\ 3M vs 1M SOFR), the market quote is the
spread $s$ such that:
\begin{equation}
  \sum_{i}\tau_i^{3M}\,F^{3M}(t_{i-1},t_i)\,D(t_i)
  =
  \sum_{j}\tau_j^{1M}\!\left[F^{1M}(t_{j-1},t_j) + s\right]\!D(t_j)
\end{equation}

Rearranging, the implied spread between the two projection curves is:
\begin{equation}
  s = \frac{\sum_i \tau_i^{3M}\,F^{3M}_i\,D_i
            - \sum_j \tau_j^{1M}\,F^{1M}_j\,D_j}
           {\sum_j \tau_j^{1M}\,D_j}
\end{equation}

Typical ranges: USD SOFR 1M/3M: $-2$ to $-5$ bp; EUR EURIBOR 1M/3M: $-3$ to $-8$ bp.

% ----------------------------------------------------------
\subsection{Cross-Currency Basis Curves}
% ----------------------------------------------------------

\sprintstatus{6}{amber!80!black}{PLANNED — Bootstrap wired Sprint 6}

\subsubsection{Covered Interest Parity}

In a frictionless market, the forward FX rate is determined by the interest
rate differential (Covered Interest Parity, CIP):
\begin{equation}
  F^{\mathrm{CIP}}(t) = S_0 \cdot \frac{D^{\mathrm{dom}}(t)}{D^{\mathrm{USD}}(t)}
\end{equation}
where $S_0$ is the FX spot rate (DOM/USD) and $D^{\mathrm{dom}}$, $D^{\mathrm{USD}}$
are the domestic and USD OIS discount factors.

Post-2008, CIP fails in practice. The observed forward $F^{\mathrm{mkt}}(t)$
deviates from $F^{\mathrm{CIP}}(t)$. The cross-currency basis $b(t)$ captures
this deviation:
\begin{equation}
  F^{\mathrm{mkt}}(t)
  = S_0 \cdot \frac{D^{\mathrm{dom}}(t)\,e^{b(t)\,t}}{D^{\mathrm{USD}}(t)}
\end{equation}

The XCCY bootstrap solves for $b(t)$ at each pillar such that a mark-to-market
XCCY basis swap (exchanging domestic OIS compounded vs USD OIS compounded) prices
to par given observed market quotes.

\subsubsection{XCCY Basis Swap Pricing}

A receive-DOM / pay-USD XCCY basis swap with basis spread $b$ on the DOM leg:
\begin{equation}
  \mathrm{NPV}
  = \underbrace{N_{\mathrm{DOM}}\left[
      \sum_i \tau_i\,\bigl(F^{\mathrm{OIS,dom}}_i + b\bigr)\,D^{\mathrm{dom}}_i
      + D^{\mathrm{dom}}(T) - 1
    \right]}_{\text{domestic leg (in DOM ccy)}}
  - \underbrace{N_{\mathrm{USD}}\,\frac{1}{S_0}\left[
      \sum_i \tau_i\,F^{\mathrm{OIS,USD}}_i\,D^{\mathrm{USD}}_i
      + D^{\mathrm{USD}}(T) - 1
    \right]}_{\text{USD leg (converted to DOM)}}
  = 0
\end{equation}

The notionals satisfy $N_{\mathrm{DOM}} = N_{\mathrm{USD}}/S_0$.

\subsubsection{XVA Implication — ColVA for Cross-Currency CSA}

When a USD trade is collateralised under a EUR CSA, the cheapest-to-deliver
(CTD) discount rate involves the XCCY basis:
\begin{equation}
  r_{\mathrm{CTD}}(t) = r^{\mathrm{\text{€STR}}}(t) + b_{\mathrm{EUR/USD}}(t)
\end{equation}

The resulting ColVA adjustment:
\begin{equation}
  \mathrm{ColVA} = -\int_0^T D(t)\,\mathbb{E}^{\mathbb{Q}}\!\left[\mathrm{EE}(t)\right]
                   \cdot \bigl(r_{\mathrm{CTD}}(t) - r_{\mathrm{OIS,USD}}(t)\bigr)\,\,\mathrm{d}t
\end{equation}

\emph{Without the XCCY curve, this calculation is incorrect for cross-currency CSA trades.}

% ----------------------------------------------------------
\subsection{Funding Curves}
% ----------------------------------------------------------

\sprintstatus{1}{teal!70!black}{COMPLETE — Manual spread UI built. XVA linkage defined.}

\subsubsection{Definition}

\begin{tcolorbox}[definition, title={Funding Curve}]
The firm's funding curve at tenor $t$ is:
\begin{equation}
  r_{\mathrm{fund}}(t) = r_{\mathrm{OIS}}(t) + \frac{s(t)}{10\,000}
\end{equation}
where $s(t)$ is the firm's credit spread above OIS in basis points.
The spread $s(t)$ is a term structure input manually by the treasury desk;
it is not bootstrapped from market instruments.
Interpolation between input tenors: log-linear, consistent with the OIS base curve.
\end{tcolorbox}

This is not bootstrapped because no directly observable market instruments
constrain the firm's funding spread term structure. The OIS bootstrap (Pass 1)
provides the risk-free base; the spread is purely additive.

\subsubsection{XVA Formulas — FVA, ColVA, MVA}

\begin{tcolorbox}[formula]
\begin{align}
  \mathrm{FVA}
  &= -\int_0^T D(t)\,\mathbb{E}^{\mathbb{Q}}\!\left[\mathrm{EE}(t) - \mathrm{ENE}(t)\right]
     \cdot \frac{s(t)}{10\,000}\;\,\mathrm{d}t
  \label{eq:fva}\\[6pt]
  \mathrm{ColVA}
  &= -\int_0^T D(t)\,\mathbb{E}^{\mathbb{Q}}\!\left[\mathrm{EE}(t)\right]
     \cdot s_{\mathrm{CTD}}(t)\;\,\mathrm{d}t
  \label{eq:colva}\\[6pt]
  \mathrm{MVA}
  &= -\int_0^T D(t)\,\mathbb{E}^{\mathbb{Q}}\!\left[\mathrm{SIMM\text{-}IM}(t)\right]
     \cdot \frac{s(t)}{10\,000}\;\,\mathrm{d}t
  \label{eq:mva}
\end{align}
\end{tcolorbox}

where:
\begin{itemize}[leftmargin=1.5cm]
  \item $D(t)$ — OIS discount factor (from base curve, Pass 1)
  \item $\mathrm{EE}(t)$ — Expected Exposure under risk-neutral measure $\mathbb{Q}$
  \item $\mathrm{ENE}(t)$ — Expected Negative Exposure
  \item $s(t)$ — firm credit spread in bp (from funding curve)
  \item $s_{\mathrm{CTD}}(t)$ — cheapest-to-deliver collateral spread
  \item $\mathrm{SIMM\text{-}IM}(t)$ — projected forward SIMM IM path (Sprint 16 input)
\end{itemize}

In discrete form (for numerical implementation, Sprint 15):
\begin{align}
  \mathrm{FVA} &\approx -\sum_{k=1}^{M} D(t_k)\,
    \bigl[\mathrm{EE}(t_k) - \mathrm{ENE}(t_k)\bigr]\cdot
    \frac{s(t_k)}{10\,000}\cdot\Delta t_k \\[4pt]
  \mathrm{MVA} &\approx -\sum_{k=1}^{M} D(t_k)\,
    \mathrm{SIMM\text{-}IM}(t_k)\cdot
    \frac{s(t_k)}{10\,000}\cdot\Delta t_k
\end{align}

\subsubsection{Sample Pre-Loaded Spreads}

Representative of a Tier 1, A-rated bank. Users replace with own institution's spreads.

\begin{center}
\small
\begin{tabular}{L{1.2cm} R{1.0cm} R{1.0cm} R{1.0cm} R{1.0cm} R{1.0cm} R{1.0cm} R{1.0cm}}
\toprule
\textbf{Tenor} & \textbf{USD} & \textbf{EUR} & \textbf{GBP} & \textbf{JPY}
              & \textbf{CHF} & \textbf{AUD} & \textbf{CAD} \\
\midrule
ON  & 12 & 14 & 10 &  5 & 10 & 15 & 12 \\
1M  & 20 & 22 & 18 &  8 & 15 & 22 & 18 \\
3M  & 28 & 30 & 25 & 12 & 22 & 30 & 25 \\
6M  & 35 & 38 & 32 & 16 & 28 & 38 & 32 \\
1Y  & 42 & 45 & 40 & 20 & 35 & 45 & 40 \\
2Y  & 52 & 55 & 50 & 25 & 42 & 55 & 50 \\
3Y  & 58 & 62 & 55 & 28 & 48 & 60 & 55 \\
5Y  & 65 & 70 & 62 & 32 & 55 & 68 & 62 \\
7Y  & 70 & 75 & 68 & 36 & 60 & —  & —  \\
10Y & 75 & 80 & 72 & 38 & 62 & 75 & 70 \\
20Y & 85 & 90 & 82 & 42 & —  & —  & —  \\
30Y & 88 & 92 & 85 & 44 & —  & —  & —  \\
\midrule
\multicolumn{8}{l}{\footnotesize All values in basis points above OIS. Blank = pillar not defined.}\\
\bottomrule
\end{tabular}
\end{center}

% ============================================================
\section{Planned Sections — Framework}
% ============================================================

\subsection{Sprint 3: OIS Bootstrap API \& History Store}

\sprintstatus{3}{amber!80!black}{PLANNED}

Will document: the live QuantLib bootstrap API endpoint, simultaneous basis curve
bootstrap implementation, three-store schema deployment, Bloomberg DAPI adapter,
FRED/ECB/BOE adapters. Mathematical content: convergence proof sketch for
Brent's method applied to OIS bootstrap; sensitivity of bootstrap result to
pillar instrument selection; condition number analysis.

\subsection{Sprint 4: Volatility Surfaces}

\sprintstatus{4}{amber!80!black}{PLANNED}

Will document SABR calibration. The SABR model (Hagan et al.\ 2002) for the
implied vol smile:
\begin{equation}
  \sigma_{\mathrm{BS}}(F,K,T) \approx
    \frac{\alpha}{(FK)^{\frac{1-\beta}{2}}}
    \cdot\frac{z}{\chi(z)}
    \cdot\left[1 + \left(\frac{(1-\beta)^2}{24}\frac{\alpha^2}{(FK)^{1-\beta}}
                          + \frac{\rho\beta\nu\alpha}{4(FK)^{\frac{1-\beta}{2}}}
                          + \frac{2-3\rho^2}{24}\nu^2\right)T\right]
\end{equation}
where $z = \frac{\nu}{\alpha}(FK)^{\frac{1-\beta}{2}}\ln(F/K)$ and
$\chi(z) = \ln\!\left(\frac{\sqrt{1-2\rho z+z^2}+z-\rho}{1-\rho}\right)$.

Parameters $(\alpha,\beta,\rho,\nu)$ calibrated per expiry slice by minimising
the weighted sum of squared vega-normalised errors across the smile.

\subsection{Sprints 7–8: Instrument Pricers}

\sprintstatus{7}{amber!80!black}{PLANNED}

Will document NPV formulas for: IRS, OIS swap, basis swap, XCCY swap, swaption
(Black and Bachelier models), cap/floor, CDS (ISDA model), CDX index.

For reference, the key pricing equations are:
\begin{align}
  \mathrm{NPV}_{\mathrm{IRS}} &=
    N\!\left[\sum_{i}\tau_i^f F_i^{\mathrm{basis}} D_i
           - K\sum_{j}\tau_j^f D_j\right]  \\[4pt]
  \mathrm{NPV}_{\mathrm{swaption}}^{\mathrm{Black}} &=
    A(T_0)\left[S(T_0,T_n)\,\Phi(d_+) - K\,\Phi(d_-)\right] \\[4pt]
  d_{\pm} &= \frac{\ln(S/K) \pm \frac{1}{2}\sigma^2 T_0}{\sigma\sqrt{T_0}}
\end{align}

\subsection{Sprint 10: VaR Engine}

\sprintstatus{10}{amber!80!black}{PLANNED}

\textbf{Historical Simulation VaR} (1-day, 99\%, 1250 scenarios):
\begin{equation}
  \mathrm{P\&L}_t = \sum_i \mathrm{NPV}_i\!\left(\mathcal{M}_t\right)
                  - \sum_i \mathrm{NPV}_i\!\left(\mathcal{M}_{t-1}\right)
\end{equation}
\begin{equation}
  \mathrm{VaR}_{0.99} = -\mathrm{Percentile}\!\left(\{\mathrm{P\&L}_t\}_{t=1}^{1250},\; 1\%\right)
\end{equation}

\textbf{Delta-Gamma VaR} (using SIMM covariance):
\begin{equation}
  \mathrm{VaR}_{\mathrm{\Delta\Gamma}} \approx z_{0.99}\sqrt{\mathbf{WS}^\T\,\Sigma\,\mathbf{WS}}
    + \frac{1}{2}\,\mathrm{tr}(\Gamma\,\Sigma)
\end{equation}

\textbf{Expected Shortfall} (for FRTB):
\begin{equation}
  \mathrm{ES}_{0.975} = -\frac{1}{n\,(1-0.975)}
    \sum_{t:\,\mathrm{P\&L}_t < \mathrm{VaR}_{0.975}} \mathrm{P\&L}_t
\end{equation}

\subsection{Sprint 15: XVA Engine}

\sprintstatus{15}{amber!80!black}{PLANNED}

Full bilateral XVA:
\begin{align}
  \mathrm{CVA} &= -\mathrm{LGD}_c\int_0^T D(t)\,\mathbb{E}^{\mathbb{Q}}[\mathrm{EE}(t)]\,\mathrm{d}\mathrm{PD}_c^{\mathbb{Q}}(t)\\
  \mathrm{DVA} &= +\mathrm{LGD}_b\int_0^T D(t)\,\mathbb{E}^{\mathbb{Q}}[\mathrm{ENE}(t)]\,\mathrm{d}\mathrm{PD}_b^{\mathbb{Q}}(t)\\
  \mathrm{KVA} &= -\int_0^T D(t)\,\mathbb{E}^{\mathbb{Q}}\!\left[\mathrm{EAD}(t)\right]
                  \cdot RW\cdot\mathrm{CoC}\;\,\mathrm{d}t
\end{align}

FVA and MVA are given by equations (\ref{eq:fva}) and (\ref{eq:mva}).

\subsection{Sprint 16: ISDA SIMM v2.6}

\sprintstatus{16}{amber!80!black}{PLANNED}

Implemented entirely in own Python. The complete aggregation:
\begin{equation}
  WS_{ik} = RW_k\cdot s_{ik}
\end{equation}
\begin{equation}
  K_b = \sqrt{\sum_i WS_{ib}^2 + \sum_{i\neq j}\rho_{ij}\,WS_{ib}\,WS_{jb}}
\end{equation}
\begin{equation}
  S_b = \max\!\left(-K_b,\;\min\!\left(K_b,\;\sum_i WS_{ib}\right)\right)
\end{equation}
\begin{equation}
  IM_{rc} = \sqrt{\sum_b K_b^2 + \sum_{b\neq c}\gamma_{bc}\,S_b\,S_c}
\end{equation}
\begin{equation}
  IM_{\mathrm{total}} = \sum_{rc} IM_{rc}
\end{equation}

% ============================================================
\section{Model Validation Guide}
% ============================================================

\subsection{OIS Bootstrap Validation Test Cases}

\begin{center}
\begin{tabular}{L{1.5cm} L{4.5cm} L{4cm} L{2.5cm}}
\toprule
\textbf{Test ID} & \textbf{Input} & \textbf{Expected Output} & \textbf{Tolerance} \\
\midrule
OIS-V001
& USD SOFR: ON=5.31\%, 1Y=4.48\%, 5Y=3.78\%, 10Y=3.74\%
& $D(1) \approx 0.9571$, $D(5) \approx 0.8289$, $D(10) \approx 0.6898$
& $10^{-4}$ on $D$ \\
\addlinespace
OIS-V002
& EUR €STR: ON=3.40\%, 1Y=2.58\%, 5Y=2.30\%
& $r(5) \approx 2.283\%$, $r(10) \approx 2.448\%$
& 0.1 bp on $r$ \\
\addlinespace
OIS-V003
& Flat curve: all tenors $= r$
& $D(t) = e^{-rt}$ exactly for all $t$
& Machine $\varepsilon$ \\
\addlinespace
OIS-V004
& Log-linear vs Cubic Spline at 3.5Y
& Forward rate difference $< 20$ bp for normal curve shapes
& Qualitative \\
\bottomrule
\end{tabular}
\end{center}

\subsection{Known Model Limitations}

\begin{center}
\begin{tabular}{L{2.5cm} L{4cm} L{2.8cm} L{3.2cm}}
\toprule
\textbf{Model} & \textbf{Limitation} & \textbf{Impact} & \textbf{Mitigation} \\
\midrule
All interpolation & Assumes continuous trading at pillar tenors; illiquid pillars ill-conditioned
& Curve poorly defined at illiquid points
& Illiquid pillars (3W, 4Y, 15Y) disabled by default \\
\addlinespace
Log-Linear Discount & Piecewise-constant forwards; unrealistic rate expectations
& Stepped forward curve; affects swaption pricing at boundaries
& Switch to Cubic or Monotone Cubic for smooth forwards \\
\addlinespace
Bootstrap accuracy $10^{-12}$ & Will not converge for inconsistent or highly non-monotone instruments
& Bootstrap fails — exception raised
& Pre-validation range checks (Sprint 3) \\
\addlinespace
Flat extrapolation & Beyond last pillar, uses last-pillar rate
& Inaccurate for ultra-long trades ($>30$Y)
& User can add longer pillar \\
\bottomrule
\end{tabular}
\end{center}

% ============================================================
\section{Regulatory Alignment}
% ============================================================

\begin{center}
\begin{tabular}{L{3cm} L{3.5cm} L{1.5cm} L{4.5cm}}
\toprule
\textbf{Requirement} & \textbf{Rijeka Module} & \textbf{Sprint} & \textbf{Standard} \\
\midrule
IM — SIMM     & SIMM engine, all 6 risk classes & 16 & ISDA SIMM v2.6 \\
IM — Schedule & Schedule IM engine              & 17 & BCBS/IOSCO WGMR \\
Variation Margin & Full revaluation MTM          &  3 & BCBS/IOSCO WGMR \\
CCR Capital   & SA-CCR EAD engine               & 14 & BCBS d424 \\
Market Risk   & FRTB ES (2520-day window)        & 10 & BCBS d457 \\
Data Lineage  & Immutable history + audit trail  &  3 & BCBS 239 \\
LCH SwapClear & FHS 1250-day replication         & 17 & LCH methodology \\
\bottomrule
\end{tabular}
\end{center}

% ============================================================
\appendix
\section{Notation Reference}
% ============================================================

\begin{center}
\begin{tabular}{L{2.5cm} L{10cm}}
\toprule
\textbf{Symbol} & \textbf{Definition} \\
\midrule
$D(t)$          & Discount factor for maturity $t$: present value of 1 unit at time $t$ \\
$r(t)$          & Continuously compounded zero rate \\
$f(t)$          & Instantaneous forward rate \\
$F(t_1,t_2)$    & Simply compounded period forward rate for $[t_1,t_2]$ \\
$S(T)$          & Par swap rate for maturity $T$ \\
$A(T)$          & Annuity (DV01 numeraire): $\sum_i \tau_i D(t_i)$ \\
$\tau(t_1,t_2)$ & Day count fraction under the relevant convention \\
$N$             & Notional \\
$K$             & Fixed rate (IRS) or strike (option) \\
$\Phi(\cdot)$   & Standard normal CDF \\
$\mathbb{E}^{\mathbb{Q}}[\cdot]$  & Expectation under risk-neutral measure $\mathbb{Q}$ \\
$\mathrm{EE}(t)$& Expected Exposure at time $t$ \\
$\mathrm{ENE}(t)$ & Expected Negative Exposure \\
$\mathrm{LGD}$  & Loss Given Default \\
$\mathrm{PD}(t)$ & Default probability (survival: $1-\mathrm{PD}$) \\
$s(t)$          & Firm credit spread at tenor $t$ (basis points) \\
$s_{ik}$        & SIMM raw sensitivity of instrument $i$ to factor $k$ \\
$RW_k$          & SIMM risk weight for factor $k$ \\
$\rho_{ij}$     & SIMM within-bucket correlation \\
$\gamma_{bc}$   & SIMM cross-bucket correlation \\
$\Sigma$        & SIMM covariance matrix: $\Sigma_{ij} = \sigma_i\sigma_j\rho_{ij}$ \\
\bottomrule
\end{tabular}
\end{center}

\section{Document Version History}

\begin{center}
\begin{tabular}{L{1.5cm} L{2cm} L{3cm} L{6cm}}
\toprule
\textbf{Version} & \textbf{Date} & \textbf{Author} & \textbf{Changes} \\
\midrule
v1.0 & 2025-01-15 & Rijeka Core Team
     & Initial release. Sprint 1 complete. Framework for Sprints 2–20. \\
v1.1 & 2025-01-15 & Rijeka Core Team
     & Rebuilt in \LaTeX. Full mathematical notation throughout.
       Sections 4.3 (Basis), 4.4 (XCCY), 4.5 (Funding) added.
       All formulas typeset. \\
\bottomrule
\end{tabular}
\end{center}

\vspace{1cm}
\begin{center}
\textcolor{mgrey}{\rule{12cm}{0.4pt}}\\[0.4cm]
{\small\color{dgrey}
\textbf{RIJEKA} — Open-Source Full Revaluation Risk System\\[0.1cm]
\textit{One system. One vocabulary. Every risk metric a derivatives book requires.}
}
\end{center}

\end{document}
